\clearpage
\begin{proofbox}
\textbf{统一记号与前提。}

设 $F_1=(-c,0),\ F_2=(c,0)$，
$P=(x_0,y_0)\in H$ 且 $y_0\ne0$，
$M=(m,0)$ 是 $\angle F_1PF_2$ 内角平分线的轴上交点。
令
\[
r_1=PF_1,\qquad r_2=PF_2,\qquad c^2=a^2+b^2.
\]
由于 $P$ 非顶点，$x_0>a$ 或 $x_0<-a$，而 $M$ 位于开线段 $F_1F_2$ 上。
以下给出\textbf{六种方法}：保留原有五种，再增加可独立推导的倒角公式法；
排序考虑入门负担、考场速度与知识迁移。

\medskip
\textbf{方法一：最低门槛【角平分线定理＋平方差；首推独立证明】。}

\begin{center}
\begin{tikzpicture}[font=\small,>=stealth]
\draw[gray!40,rounded corners] (0,0) rectangle (14.2,9);
\node[font=\small\bfseries] at (7.1,8.35) {方法一：两个直角三角形，共用同一条高};

\begin{scope}[shift={(4.4,3.2)},scale=.6]

\coordinate (L) at (-5,0); \coordinate (R) at (5,0);
\coordinate (P) at (3.75,3); \coordinate (M) at (2.4,0);
\draw[blue!75!black,thick] (L)--(P);
\draw[orange!85!black,thick] (P)--(R);
\draw[gray] (L)--(R);
\draw[magenta!80!black,very thick] (M)--(P);
\foreach \q in {L,R,P,M} \fill (\q) circle (3pt);
\node[below left] at (L) {$F_1$}; \node[below right] at (R) {$F_2$};
\node[above left] at (P) {$P$}; \node[below left] at (M) {$M$};

\coordinate (H) at (3.75,0);
\draw[dashed,gray] (P)--(H);
\draw (3.45,0)--(3.45,.3)--(3.75,.3);
\fill (H) circle(3pt);\node[below] at (H) {$H$};
\draw[blue!75!black,<->] (-5,-1.8)--(3.75,-1.8) node[midway,fill=white] {$x_0+c$};
\draw[orange!85!black,<->] (3.75,-2.7)--(5,-2.7);
\node[orange!85!black,anchor=west] at (5.4,-2.7) {$|x_0-c|$};
\node[blue!75!black,above left] at (-.6,1.5) {$r_1$};
\node[orange!85!black,right] at (5.2,1.7) {$r_2$};
\end{scope}
\node[align=left,anchor=west] at (9,5.5) {$r_1^2=(x_0+c)^2+y_0^2$\\[12pt]$r_2^2=(x_0-c)^2+y_0^2$\\[12pt]相减，公共的高消去：\\[6pt]$r_1^2-r_2^2=4cx_0$};
\node[anchor=west] at (9,2.3) {$H=(x_0,0),\quad PH=|y_0|$};
\node at (7.1,.65) {图示 $0<x_0<c$；水平长度取绝对值，平方公式对左右两支均成立。};

\end{tikzpicture}
\end{center}

\textbf{第1步：角平分线定理。}
\[
\frac{c+m}{c-m}=\frac{r_1}{r_2}
\quad\Longrightarrow\quad
r_2(c+m)=r_1(c-m),
\]
于是
\[
\boxed{m=c\frac{r_1-r_2}{r_1+r_2}}.\tag{1}
\]
\textbf{第2步：焦距平方相减。}
\[
\begin{aligned}
r_1^2-r_2^2
&=(x_0+c)^2+y_0^2-(x_0-c)^2-y_0^2\\
&=4cx_0.
\end{aligned}
\]
即
\[
\boxed{x_0=\frac{(r_1-r_2)(r_1+r_2)}{4c}}.\tag{2}
\]
\textbf{第3步：乘起来，使用双曲线第一定义。}
\[
\begin{aligned}
mx_0
&=\frac{c(r_1-r_2)}{r_1+r_2}
  \cdot\frac{(r_1-r_2)(r_1+r_2)}{4c}\\
&=\frac{(r_1-r_2)^2}{4}
=\frac{(2a)^2}{4}=\boxed{a^2}.
\end{aligned}
\]
\textbf{优势：}完全不必区分左右支；不背焦半径；只用角平分线定理、
两点距离与平方差公式。\textbf{关键切入：}两个焦距是“和、差”，而 $x_0$ 是“平方差”。

\medskip
\end{proofbox}
\clearpage
\begin{proofbox}
\textbf{方法二：最快代数【角平分线定理＋焦半径】。}

先沿用式 (1)。令 $e=c/a>1$，记
$\sigma=\operatorname{sgn}(x_0)\in\{1,-1\}$。
双曲线左右两支的焦半径可以统一写成
\[
\boxed{r_1=e|x_0|+\sigma a,\qquad
r_2=e|x_0|-\sigma a}.
\]
核对：右支 $\sigma=1$ 时 $r_1=ex_0+a,\ r_2=ex_0-a$；
左支 $\sigma=-1$ 时 $r_1=-ex_0-a,\ r_2=-ex_0+a$。
因此
\[
\begin{aligned}
m
&=c\frac{2\sigma a}{2e|x_0|}
=\frac{ca\sigma}{e|x_0|}
=\frac{a^2\sigma}{|x_0|}
=\boxed{\frac{a^2}{x_0}}.
\end{aligned}
\]
\textbf{优势：}已熟悉焦半径时只需代入一次。
\textbf{代价：}若左右焦点匹配或左右支的符号记错，很容易出错。

\medskip
\end{proofbox}
\clearpage
\begin{proofbox}
\textbf{方法三：最快识别【已知切线平分焦点三角形的内角】。}

若已掌握双曲线的\textbf{切线--焦点角平分线性质}，
$PM$ 就是 $H$ 在 $P$ 处的切线。
切线方程为
\[
\frac{x_0x}{a^2}-\frac{y_0y}{b^2}=1.
\]
令 $y=0$，得到
\[
\frac{x_0m}{a^2}=1
\quad\Longrightarrow\quad
\boxed{m=\frac{a^2}{x_0}}.
\]
\textbf{优势：}看出切线后，求值几乎一步完成。
\textbf{注意：}本方法\emph{调用了已知的切线--角平分线性质}，
不能再把这一步当作对该性质的独立证明；若还没掌握该性质，用方法一建立它。

\medskip
\end{proofbox}
\clearpage
\begin{proofbox}
\textbf{方法四：几何深化【内切圆切线长＋角平分线定理】。}

\begin{center}
\begin{tikzpicture}[font=\small,>=stealth]
\draw[gray!40,rounded corners] (0,0) rectangle (14.2,8.5);
\node[font=\small\bfseries] at (7.1,7.85) {方法四：先看三角形的等切线段，再代入焦距差};

\begin{scope}[shift={(4.3,3.1)},scale=.68]

\coordinate (L) at (-5,0); \coordinate (R) at (5,0);
\coordinate (P) at (3.75,3); \coordinate (M) at (2.4,0);
\draw[blue!75!black,thick] (L)--(P);
\draw[orange!85!black,thick] (P)--(R);
\draw[gray] (L)--(R);
\draw[magenta!80!black,very thick] (M)--(P);
\foreach \q in {L,R,P,M} \fill (\q) circle (3pt);
\node[below left] at (L) {$F_1$}; \node[below right] at (R) {$F_2$};
\node[above left] at (P) {$P$}; \node[below left] at (M) {$M$};

\coordinate (I) at (3,1.333333);\coordinate (T) at (3,0);
\coordinate (U) at (2.567568,2.594595);\coordinate (V) at (4.230769,1.846154);
\draw[teal!80!black,thick] (I) circle(1.333333);
\draw[teal!80!black,dashed] (I)--(T);
\draw[blue!75!black,very thick] (L)--(U);\draw[blue!75!black,very thick] (L)--(T);
\draw[orange!85!black,very thick] (R)--(V);\draw[orange!85!black,very thick] (R)--(T);
\draw[purple,very thick] (U)--(P)--(V);
\foreach \q in {I,T,U,V} \fill (\q) circle(3pt);
\node[left] at (I) {$I$}; \draw (T)--(3,-1) node[below] {$T$};
\node[above left] at (U) {$U$};\node[right] at (V) {$V$};
\end{scope}
\node[anchor=west,align=left] at (9,5.5) {$F_1U=F_1T$（蓝）\\[8pt]$F_2V=F_2T$（橙）\\[8pt]$PU=PV=\lambda$（紫）};
\node[align=center] at (7.1,1) {$r_1-r_2=F_1T-F_2T$，公共的 $\lambda$ 抵消。\\右支差为 $2a$，所以 $T=(a,0)$；而 $M=(a^2/x_0,0)$。};

\end{tikzpicture}
\end{center}

作 $\triangle PF_1F_2$ 的内切圆，切 $F_1F_2$ 于 $T(t,0)$，
切 $PF_1$、$PF_2$ 于 $U,V$。令 $PU=PV=\lambda>0$。
同一点的两条切线段等长，所以
\[
r_1=F_1T+\lambda,\qquad
r_2=F_2T+\lambda.
\]
令 $\sigma=\operatorname{sgn}(x_0)$，故
\[
F_1T-F_2T=r_1-r_2=2\sigma a,
\qquad F_1T+F_2T=2c,
\]
\[
F_1T=c+\sigma a,\qquad
F_2T=c-\sigma a,
\qquad\boxed{T=(\sigma a,0)}.
\]
这先证明了\textbf{内切圆与焦点轴的切点始终为对应分支顶点}。
并且
\[
r_1=\lambda+c+\sigma a,\qquad
r_2=\lambda+c-\sigma a,
\]
由距离平方差
\[
4cx_0=r_1^2-r_2^2
=(2\sigma a)\cdot2(\lambda+c),
\]
得
\[
\lambda+c=\frac{c|x_0|}{a}.
\]
最后用角平分线定理的式 (1)：
\[
m=c\frac{r_1-r_2}{r_1+r_2}
=c\frac{2\sigma a}{2(\lambda+c)}
=\boxed{\frac{a^2}{x_0}}.
\]
\textbf{优势：}揭示定切点 $T=\pm a$，几何画面非常有记忆点。
\textbf{代价：}为了进一步求 $M$，还需连接焦距平方差或其他长度关系；
不如方法一直接。切记 $T$ 与 $M$ 通常不重合。

\medskip
\end{proofbox}
\clearpage
\begin{proofbox}
\textbf{方法五：半角三角函数【单位方向向量的内角平分线】。}

\begin{center}
\begin{tikzpicture}[font=\small,>=stealth]
\draw[gray!40,rounded corners] (0,0) rectangle (14.2,8);
\node[font=\small\bfseries] at (7.1,7.35) {方法五：先把两条射线各取一段单位长度};

\begin{scope}[shift={(5,5.6)},scale=2]
\coordinate (P) at (0,0); \coordinate (E) at (-.945946,-.324324);
\coordinate (G) at (.384615,-.923077);\coordinate (S) at (-.561331,-1.247401);
\draw[gray] (0,0)--(-1.55,-.53143);\draw[gray] (0,0)--(.61538,-1.47692);
\draw[blue!75!black,very thick,->] (P)--(E) node[midway,above left] {$\boldsymbol u_1$};
\draw[orange!85!black,very thick,->] (P)--(G) node[midway,right] {$\boldsymbol u_2$};
\draw[dashed] (E)--(S)--(G);
\draw[magenta!80!black,very thick,->] (P)--(S);
\foreach \q in {P,E,G,S} \fill (\q) circle(.8pt);
\node[above] at (P) {$P$};\node[left] at (E) {$E$};\node[right] at (G) {$G$};
\node[below] at (S) {$S$};
\end{scope}
\node[anchor=west,align=left] at (9,4.8) {$PE=PG=1$\\[8pt]平行四边形成为菱形\\[8pt]对角线 $PS$ 平分顶角};
\node[align=center] at (7.1,1) {$\overrightarrow{PS}=\boldsymbol u_1+\boldsymbol u_2$ 沿内角平分线。\\$PE$ 指向 $F_1$，$PG$ 指向 $F_2$；单位化后才能直接相加。};

\end{tikzpicture}
\end{center}

设两条从 $P$ 指向焦点的\textbf{单位方向向量}为
\[
\boldsymbol{u}_1=\frac{(-c-x_0,-y_0)}{r_1},
\qquad
\boldsymbol{u}_2=\frac{(c-x_0,-y_0)}{r_2}.
\]
若顶角 $\gamma=\angle F_1PF_2\in(0,\pi)$，
两个单位向量相加得到
\[
\boldsymbol{u}_1+\boldsymbol{u}_2
=2\cos\frac{\gamma}{2}\,\boldsymbol{w},
\]
其中 $\boldsymbol{w}$ 是\textbf{内角平分线的单位方向}。
这一等式来自\emph{半角余弦公式}；因此
$\boldsymbol{v}=\boldsymbol{u}_1+\boldsymbol{u}_2$ 就沿 $PM$。

设 $M=P+k\boldsymbol v$。利用 $M$ 的纵坐标为 $0$，
\[
0=y_0-k y_0\Bigl(\frac1{r_1}+\frac1{r_2}\Bigr)
\quad\Longrightarrow\quad
k=\frac{r_1r_2}{r_1+r_2}.
\]
代入横坐标：
\[
\begin{aligned}
m
&=x_0+k\Bigl(-\frac{x_0+c}{r_1}+\frac{c-x_0}{r_2}\Bigr)\\
&=x_0+\frac{-(x_0+c)r_2+(c-x_0)r_1}{r_1+r_2}\\
&=c\frac{r_1-r_2}{r_1+r_2}.
\end{aligned}
\]
再用 $r_1^2-r_2^2=4cx_0$ 与 $|r_1-r_2|=2a$，
便得 $mx_0=a^2$。
\textbf{优势：}直接由\emph{半角方向}建立直线，无需预先调用角平分线定理。
\textbf{代价：}向量/三角运算更多，不适合追求最低认知负担的首解。

\medskip
\end{proofbox}

\clearpage
\begin{proofbox}
\textbf{方法六：倒角公式法【两直线夹角的正切公式；独立证明】。}

\textbf{目标：}不预先知道“$PM$ 是切线”，也不用角平分线定理；
直接利用 $\angle F_1PM=\angle MPF_2$，先求角平分线的轴上交点 $M$。

% Same Cartesian scale in both panels; a=3, b=4, c=5, P=(15/4,3).
\begin{center}
\begin{tikzpicture}[font=\small,>=stealth]
\draw[gray!40,rounded corners] (0,0) rectangle (14.2,13.5);
\node[font=\small\bfseries] at (7.1,12.8)
 {倒角公式：先看两个等角，再看两个交点};
\node[anchor=west] at (.6,11.8) {① 只看普通三角形：紫线把顶角分成两个 $\alpha$};
\begin{scope}[shift={(4.3,8.6)},scale=.68]
 \coordinate (L) at (-5,0); \coordinate (R) at (5,0);
 \coordinate (P) at (3.75,3); \coordinate (M) at (2.4,0);
 \fill[blue!5] (L)--(P)--(M)--cycle;
 \fill[orange!8] (M)--(P)--(R)--cycle;
 \draw[gray] (L)--(R);
 \draw[blue!75!black,thick] (L)--(P) node[midway,above left] {$k_1$};
 \draw[orange!85!black,thick] (P)--(R) node[midway,right] {$k_2$};
 \draw[magenta!80!black,very thick] (M)--(P);
 \draw[magenta!80!black] (P) ++(-161.075356:1.1)
   arc (-161.075356:-114.227745:1.1);
 \draw[magenta!80!black] (P) ++(-114.227745:1.4)
   arc (-114.227745:-67.380135:1.4);
 \node at (2.05,1.9) {$\alpha$};
 \node at (3.85,1.05) {$\alpha$};
 \foreach \q in {L,R,P,M} \fill (\q) circle(2.6pt);
 \node[below left] at (L) {$F_1$}; \node[below right] at (R) {$F_2$};
 \node[above] at (P) {$P$}; \node[below left] at (M) {$M$};
\end{scope}
\node[anchor=west,align=left] at (9.35,10.1)
 {$k_1,k_2$：两腰的斜率\\[8pt]$k$：紫线 $PM$ 的斜率};
\node[anchor=west,align=left] at (9.35,8.4)
 {两边都在算\\同一个 $\tan\alpha$};
\draw[gray!25] (.5,7.5)--(13.7,7.5);
\node[anchor=west] at (.6,6.9) {② 放回本题：正切等式还会给出外角平分线};
\begin{scope}[shift={(4.3,3.8)},scale=.68]
 \coordinate (L) at (-5,0); \coordinate (R) at (5,0);
 \coordinate (P) at (3.75,3); \coordinate (M) at (2.4,0);
 \coordinate (N) at (10.416666667,0);
 \draw[->,gray] (-5.6,0)--(12.2,0) node[right] {$x$};
 \draw[blue!75!black,thick] (L)--(P);
 \draw[blue!75!black,densely dotted] (P)--(6.2,3.84);
 \draw[orange!85!black,thick] (P)--(R);
 \draw[magenta!80!black,very thick] (M)--(P);
 \draw[teal!75!black,thick,dashed] (P)--(N);
 % The outward ray PN is perpendicular to the internal ray PM.
 \draw[gray] (P) ++(-114.227745:.42)
   -- ++(-24.227745:.42) -- ++(65.772255:.42);
 \draw[teal!75!black] (P) ++(-67.380135:1.05)
   arc (-67.380135:-24.227745:1.05);
 \draw[teal!75!black] (P) ++(-24.227745:1.4)
   arc (-24.227745:18.924644:1.4);
 \node[teal!75!black,fill=white,inner sep=1pt] at (5.05,1.7) {$\beta$};
 \node[teal!75!black,fill=white,inner sep=1pt] at (6,2.8) {$\beta$};
 \foreach \q in {L,R,P,M,N} \fill (\q) circle(2.6pt);
 \node[below right] at (L) {$F_1(-c,0)$};
 \node[below right] at (R) {$F_2(c,0)$};
 \node[above left] at (P) {$P$};
 \node[below left] at (M) {$M$}; \node[below] at (N) {$N$};
 \node[teal!75!black,above right] at (7.8,1.1) {外角平分线};
\end{scope}
\node[magenta!80!black,align=center] at (4.1,1.8)
 {$M\left(a^2/x_0,0\right)$\\[6pt]底边内部：$-c<m<c$，保留};
\node[teal!75!black,align=center] at (10.5,1.8)
 {$N\left(c^2x_0/a^2,0\right)$\\[6pt]底边外部：$|x_N|>c$，舍去};
\node at (7.1,.55) {图取例6的数据；两幅几何图同尺度，左支情形左右反射。};
\end{tikzpicture}
\end{center}

\textbf{看图读式：}上图先固定两个相等半角；下图把另一根命名为 $N$，
用底边内部／外部区分内、外角平分线。下文仍以 $m$ 表示待求的候选轴截距。
\end{proofbox}
\clearpage
\begin{proofbox}
\textbf{方法六（续）：由等角列式，再用位置选根。}

\textbf{第1步：设三条直线的斜率。}
在斜率均存在的通常情形（$x_0\ne\pm c$、$m\ne x_0$）下，记
\[
\boxed{k_1=\frac{y_0}{x_0+c}},\qquad
\boxed{k_2=\frac{y_0}{x_0-c}},\qquad
\boxed{k=\frac{y_0}{x_0-m}}.
\]
这里分子为 $y_0$ 而不是 $-y_0$：例如
$k_1=(-y_0)/(-c-x_0)=y_0/(x_0+c)$。

\textbf{第2步：写出相邻两角的正切。}
由 $\tan(\alpha-\beta)=\dfrac{\tan\alpha-\tan\beta}
{1+\tan\alpha\tan\beta}$，角平分条件导出
\[
\boxed{\frac{k-k_1}{1+kk_1}=\frac{k_2-k}{1+k_2k}}.\tag{6.1}
\]
注意：式 (6.1) 是\textbf{有向夹角正切}的等式，
单凭正切相等可能还包含外角平分线；必须在最后检查 $-c<m<c$。

\textbf{第3步：代入并消去斜率。}
\[
\frac{k-k_1}{1+kk_1}
=\frac{y_0(c+m)}{(x_0-m)(x_0+c)+y_0^2},
\]
\[
\frac{k_2-k}{1+k_2k}
=\frac{y_0(c-m)}{(x_0-m)(x_0-c)+y_0^2}.
\]
记
\[
D_1=(x_0-m)(x_0+c)+y_0^2,\qquad
D_2=(x_0-m)(x_0-c)+y_0^2.
\]
因为 $P$ 不是顶点，$y_0\ne0$；对角平分线，两半角均小于 $90^\circ$，
相应分母不为零。因此交叉相乘得到
\[
(c+m)D_2-(c-m)D_1=0.\tag{6.2}
\]

\textbf{第4步：只用双曲线方程，因式分解。}
由 $P\in H$ 与 $c^2=a^2+b^2$，
\[
\boxed{y_0^2=\frac{b^2}{a^2}(x_0^2-a^2)},\qquad b^2=c^2-a^2.
\]
代入式 (6.2) 并整理（可直接按 $m$ 收集同类项）：
\[
\begin{aligned}
0&=(c+m)D_2-(c-m)D_1\\
 &=2\bigl[c^2(m-x_0)+m(x_0^2-mx_0+y_0^2)\bigr]\\
 &=\frac{2}{a^2}\underbrace{(mx_0-a^2)}_{\text{内角平分线}}
                   \underbrace{(c^2x_0-a^2m)}_{\text{外角平分线}}.
\end{aligned}
\]
因此
\[
\boxed{m=\frac{a^2}{x_0}}\quad\text{或}\quad
\boxed{m=\frac{c^2x_0}{a^2}}.
\]

\textbf{第5步：用“内角平分线交底边内部”筛选。}
由双曲线非顶点条件 $|x_0|>a$、$c>a$，
\[
\left|\frac{a^2}{x_0}\right|<a<c,
\qquad
\left|\frac{c^2x_0}{a^2}\right|>\frac{c^2}{a}>c.
\]
故只有第一个根落在 $-c<m<c$ 内：
\[
\boxed{M=\left(\frac{a^2}{x_0},0\right)},\qquad
\boxed{x_0m=a^2}.
\]
另一个根对应图中的 $N$，位于底边延长线上，应舍去。

\textbf{补充：为什么第二个根是外角平分线？}
两根分别对应直线斜率
\[
 k_{\mathrm{in}}=\frac{y_0}{x_0-a^2/x_0}
 =\frac{b^2x_0}{a^2y_0},\qquad
 k_{\mathrm{out}}=\frac{y_0}{x_0-c^2x_0/a^2}
 =-\frac{a^2y_0}{b^2x_0}.
\]
于是 $\boxed{k_{\mathrm{in}}k_{\mathrm{out}}=-1}$：
两条角平分线互相垂直，正是内、外角平分线的几何特征。

\medskip
\textbf{特殊情形补全（某条焦点连线竖直时）。}
当 $x_0=c$ 或 $x_0=-c$ 时，$k_2$ 或 $k_1$ 不存在，
不能使用式 (6.1) 的斜率形式。但正切差角公式可改用
\textbf{二维向量的行列式除以点积}，不需任何直线斜率：
\[
\boldsymbol u=(-c-x_0,-y_0)=\overrightarrow{PF_1},\quad
\boldsymbol v=(c-x_0,-y_0)=\overrightarrow{PF_2},\quad
\boldsymbol w=(m-x_0,-y_0)=\overrightarrow{PM}.
\]
对应的两个半角正切分别为
\[
\frac{\det(\boldsymbol u,\boldsymbol w)}
     {\boldsymbol u\cdot\boldsymbol w}
=\frac{y_0(c+m)}{D_1},\qquad
\frac{\det(\boldsymbol w,\boldsymbol v)}
     {\boldsymbol w\cdot\boldsymbol v}
=\frac{y_0(c-m)}{D_2}.
\]
这恰好重复式 (6.2)，且对 $x_0=\pm c$ 仍有效，
因此本方法覆盖左右两支的全部非顶点点。

\textbf{评价：}这是一种真正独立于角平分线定理、焦半径公式、
切线性质的解析证明；它揭示了\textbf{两个代数根分别对应内、外角平分线}。
代价是分母和符号需要仔细处理，因此适合作为补充证明，而非最低门槛首解。
\end{proofbox}

\clearpage
\begin{proofbox}
\textbf{补充验算：从 $x_0m=a^2$ 独立推出切线方程。}

因 $y_0\ne0$，经过 $P(x_0,y_0)$、$M(a^2/x_0,0)$ 的直线满足
\[
y(x_0^2-a^2)=y_0(xx_0-a^2).
\]
又因 $P\in H$，
\[
\frac{x_0^2}{a^2}-\frac{y_0^2}{b^2}=1
\quad\Longrightarrow\quad
x_0^2-a^2=\frac{a^2y_0^2}{b^2}.
\]
代回并除以 $a^2y_0$：
\[
\frac{y_0y}{b^2}=\frac{x_0x}{a^2}-1,
\]
即
\[
\boxed{PM:\frac{x_0x}{a^2}-\frac{y_0y}{b^2}=1}.
\]
这是双曲线在 $P$ 处的标准切线方程，
因而\textbf{用方法一不仅证明截距，还独立证明了内角平分线即切线}。

\medskip
\textbf{选法建议：}需要\emph{最低门槛独立证明}用方法一；
已熟悉焦半径且追求代数速度用方法二；
切线性质已是已知结论时用方法三；
要记住内切圆固定切点用方法四；
研究内角平分线的半角方向用方法五；
要用\textbf{两直线夹角正切公式独立代数证明}，或同时研究内、外角平分线，
用方法六。
\end{proofbox}
